\documentclass[11pt]{article}
\usepackage[margin=1in]{geometry}
\usepackage[T1]{fontenc}
\usepackage{lmodern}
\usepackage{amsmath,amssymb,amsthm}
\usepackage{hyperref}
\hypersetup{colorlinks=true,linkcolor=blue,citecolor=blue,urlcolor=blue}

\newtheorem{theorem}{Theorem}
\newtheorem{corollary}{Corollary}
\newcommand{\R}{\mathbb{R}}
\newcommand{\F}{\widehat{f}}
\newcommand{\PhiA}{\Phi_{a,\sigma}}

\title{Polynomial--Gaussian Test Families and Finite Weil Diagnostics}
\author{Rajeev Kumar (Quantovibe Research Group)}
\date{Version 0.2.0, 2026}

\begin{document}
\maketitle

\begin{abstract}
We record a two-parameter Gaussian family and its polynomial extensions,
which enforce prescribed complex Fourier zeros while retaining real-even
profiles. We derive inverse transforms and autocorrelations, then formulate
an unconditional $L^2$ Gram matrix and a finite-cutoff Weil-form diagnostic.
The latter has no certified prime-tail bound. No result about the zeros of
the Riemann zeta function is asserted here.
\end{abstract}

\section{Convention and family}

Our Fourier transform and inverse are
\begin{equation}\label{eq:convention}
  \widehat{f}(\xi)=\int_{\R} f(x)e^{-i\xi x}\,dx,
  \qquad
  f(x)=\frac{1}{2\pi}\int_{\R}\widehat{f}(\xi)e^{i\xi x}\,d\xi.
\end{equation}
Fix a center $a\geq 0$ and width $\sigma>0$. Define
\begin{equation}\label{eq:profile}
  F_{a,\sigma}(\xi)
  =\left(\xi^2+\frac14\right)
    \left(e^{-(\xi-a)^2/(2\sigma^2)}
         +e^{-(\xi+a)^2/(2\sigma^2)}\right),
  \qquad \widehat{f_{a,\sigma}}=F_{a,\sigma}.
\end{equation}
The right-hand side is entire. Its restriction to $\R$ is real, even, and
Schwartz, and $F_{a,\sigma}(i/2)=F_{a,\sigma}(-i/2)=0$. Thus its inverse
transform $f_{a,\sigma}$ is real, even, and Schwartz. The complex zeros
follow from the factor $\xi^2+1/4$; they do not imply a claim about zeta
zeros.

\begin{theorem}[Inverse transform]\label{thm:inverse}
For real $x$,
\begin{align}\label{eq:inverse}
f_{a,\sigma}(x)
={}&\sqrt{\frac{2}{\pi}}\,\sigma e^{-\sigma^2x^2/2}
\Bigl[
  \left(a^2+\sigma^2+\frac14-\sigma^4x^2\right)\cos(ax)
  -2a\sigma^2x\sin(ax)
\Bigr].
\end{align}
\end{theorem}

\begin{proof}
The inverse transform of $e^{-(\xi-b)^2/(2\sigma^2)}$ under
\eqref{eq:convention} is
$\sigma e^{-\sigma^2x^2/2}e^{ibx}/\sqrt{2\pi}$.
Therefore the inverse transform of the Gaussian pair in
\eqref{eq:profile} is
$g(x)=\sqrt{2/\pi}\,\sigma e^{-\sigma^2x^2/2}\cos(ax)$.
Multiplication by $\xi^2+1/4$ corresponds to applying
$-\partial_x^2+1/4$ to $g$. Differentiating gives
\eqref{eq:inverse}.
\end{proof}

\section{Autocorrelation}

Put $\widetilde f(x)=f(-x)$ and
\begin{equation}\label{eq:autocorrelation}
  \PhiA(x)=(f_{a,\sigma}*\widetilde f_{a,\sigma})(x)
  =\int_{\R}f_{a,\sigma}(t)f_{a,\sigma}(t-x)\,dt.
\end{equation}
As $f_{a,\sigma}$ is real and even,
$\widehat{\PhiA}(\xi)=F_{a,\sigma}(\xi)^2$.

For a compact closed form, define
\begin{align}\label{eq:polynomial}
  \lambda_b(x)&=ib-\frac{\sigma^2x}{2},\\
  P_\sigma(z)&=z^4-\left(3\sigma^2+\frac12\right)z^2
       +\frac{3\sigma^4}{4}+\frac{\sigma^2}{4}+\frac1{16}.
\end{align}

\begin{theorem}[Exact autocorrelation]\label{thm:auto}
For every real $x$,
\begin{equation}\label{eq:closed}
  \PhiA(x)=\frac{\sigma}{\sqrt\pi}e^{-\sigma^2x^2/4}
  \left[
    \operatorname{Re}\!\left(e^{iax}P_\sigma(\lambda_a(x))\right)
    +e^{-a^2/\sigma^2}P_\sigma(\lambda_0(x))
  \right].
\end{equation}
Equivalently, if
\begin{equation}\label{eq:base}
  H(x)=\frac{\sigma}{\sqrt\pi}e^{-\sigma^2x^2/4}
             \left[\cos(ax)+e^{-a^2/\sigma^2}\right],
\end{equation}
then $\PhiA=(-\partial_x^2+1/4)^2H$.
\end{theorem}

\begin{proof}
Squaring the two Gaussian terms in \eqref{eq:profile} gives
\begin{align*}
\left(e^{-(\xi-a)^2/(2\sigma^2)}
     +e^{-(\xi+a)^2/(2\sigma^2)}\right)^2
={}&e^{-(\xi-a)^2/\sigma^2}
    +e^{-(\xi+a)^2/\sigma^2}
    +2e^{-a^2/\sigma^2}e^{-\xi^2/\sigma^2}.
\end{align*}
The inverse transform of this sum is $H$ in \eqref{eq:base}.
The remaining factor $(\xi^2+1/4)^2$ proves the differential form.

To evaluate it, let
$q_b(x)=e^{-\sigma^2x^2/4+ibx}$. Then
$q_b''=(\lambda_b^2-\sigma^2/2)q_b$ and
$q_b^{(4)}=(\lambda_b^4-3\sigma^2\lambda_b^2
                  +3\sigma^4/4)q_b$.
Since $(-\partial_x^2+1/4)^2
=\partial_x^4-\tfrac12\partial_x^2+\tfrac1{16}$,
its action on $q_b$ is $P_\sigma(\lambda_b)q_b$.
Pairing the conjugate terms $b=a$ and $b=-a$ yields
\eqref{eq:closed}.
\end{proof}

\begin{corollary}[Normalization checks]\label{cor:checks}
The total autocorrelation mass and its value at the origin are
\begin{align}\label{eq:mass}
  \int_{\R}\PhiA(x)\,dx
  &=F_{a,\sigma}(0)^2
   =\frac14 e^{-a^2/\sigma^2},\\
  \PhiA(0)
  &=\|f_{a,\sigma}\|_{L^2(\R)}^2\notag\\
  &=\frac{\sigma}{\sqrt\pi}
    \left[
      a^4+\left(3\sigma^2+\frac12\right)a^2
      +C_\sigma\left(1+e^{-a^2/\sigma^2}\right)
    \right],\label{eq:at-zero}
\end{align}
where $C_\sigma=3\sigma^4/4+\sigma^2/4+1/16$.
Moreover $|\PhiA(x)|\leq\PhiA(0)$ for all real $x$.
\end{corollary}

\begin{proof}
The first identity is the Fourier transform of \eqref{eq:autocorrelation}
at zero. The second follows from \eqref{eq:autocorrelation} at zero
and \eqref{eq:closed}; it also follows by Plancherel from
$(2\pi)^{-1}\int_{\R}F_{a,\sigma}(\xi)^2\,d\xi$.
The final bound is Cauchy--Schwarz applied to the two translated copies
of $f_{a,\sigma}$ in \eqref{eq:autocorrelation}.
\end{proof}

\section{Polynomial--Gaussian families}

Let $P(y)=\sum_{m=0}^{M}c_m y^m$ be a polynomial with real
coefficients. Replacing the factor $\xi^2+1/4$ in
\eqref{eq:profile} by $P(\xi^2)$ defines
\begin{equation}\label{eq:poly-family}
 F_{P,a,\sigma}(\xi)=P(\xi^2)
 \left(e^{-(\xi-a)^2/(2\sigma^2)}
       +e^{-(\xi+a)^2/(2\sigma^2)}\right).
\end{equation}
For a requested complex zero $z$, a real factor can be formed as
$Q_z(y)=y-z^2$ when $z^2$ is real, and otherwise as
$Q_z(y)=(y-z^2)(y-\overline z^{\,2})$.
Products of such factors impose the zero orbit
$\{\pm z,\pm\overline z\}$ as needed while keeping
$F_{P,a,\sigma}$ real and even on the real axis. Repeated generators
give repeated factors. The zero orbit containing $\pm i/2$ must be present
when the admissibility condition $F_{P,a,\sigma}(\pm i/2)=0$ is required
for the Weil formulas below.

\begin{theorem}[Polynomial inverse and autocorrelation]\label{thm:poly}
Let $g(x)=\sqrt{2/\pi}\,\sigma e^{-\sigma^2x^2/2}\cos(ax)$ and let
$H$ be defined by \eqref{eq:base}. Then
\begin{align}\label{eq:poly-operators}
 f_{P,a,\sigma}&=P(-\partial_x^2)g,\\
 \Phi_{P,a,\sigma}&=P(-\partial_x^2)^2H.\notag
\end{align}
Both functions are real, even, and Schwartz.
\end{theorem}

\begin{proof}
Multiplication by $\xi^{2m}$ under \eqref{eq:convention}
corresponds to $(-\partial_x^2)^m$. Apply this term by term to
the Gaussian pair and then to its squared Fourier profile.
Real coefficients and even powers preserve reality and evenness.
\end{proof}

The implementation evaluates the derivatives in
\eqref{eq:poly-operators} with the probabilists' Hermite polynomials
$\mathrm{He}_n(t)=(-1)^n e^{t^2/2}
(\mathrm{d}/\mathrm{d}t)^n e^{-t^2/2}$.
In particular, if $P(y)=\sum_m c_m y^m$, then
\begin{equation}\label{eq:hermite-inverse}
 f_{P,a,\sigma}(x)=\sqrt{\frac2\pi}\,\sigma e^{-\sigma^2x^2/2}
 \operatorname{Re}\!\left[
 e^{iax}\sum_{m=0}^{M}c_m(-1)^m\sigma^{2m}
 \mathrm{He}_{2m}\!\left(\sigma x-\frac{ia}{\sigma}\right)
 \right].
\end{equation}
For autocorrelation, expand $P(y)^2=\sum_m d_m y^m$ and apply
the same identity to the three Gaussian terms in the proof of
Theorem~\ref{thm:auto}, with Gaussian scale $\sigma/\sqrt2$.
These expressions make arbitrary finite zero sets computable
without numerical differentiation.

\section{Gram matrices and finite Weil diagnostics}

Consider a finite collection of real-even families
$f_1,\ldots,f_r$ with Fourier profiles $F_1,\ldots,F_r$.
Their cross-correlations are
\begin{equation}\label{eq:cross}
 \Phi_{ij}(x)=\int_{\R}f_i(t)f_j(t-x)\,dt
 =\frac1{2\pi}\int_{\R}F_i(\xi)F_j(\xi)e^{i\xi x}\,d\xi.
\end{equation}
At zero lag the $L^2$ Gram matrix is
\begin{equation}\label{eq:l2gram}
 G_{ij}=\Phi_{ij}(0)
 =\frac1{2\pi}\int_{\R}F_i(\xi)F_j(\xi)\,d\xi.
\end{equation}
\begin{theorem}[Unconditional Gram positivity]\label{thm:gram}
The real symmetric matrix $G$ in \eqref{eq:l2gram} is positive
semidefinite. It is positive definite if and only if the functions
$f_1,\ldots,f_r$ are linearly independent in $L^2(\R)$.
\end{theorem}
\begin{proof}
For any real vector $v$,
$v^\mathsf{T}Gv=\int_{\R}(\sum_i v_i f_i(x))^2\,dx\geq0$.
Equality holds precisely for a linear dependence.
\end{proof}
Eigenvalues and condition numbers of $G$ therefore diagnose
linear dependence and numerical conditioning. This positivity
is an ordinary $L^2$ fact; it is not a test of the Riemann Hypothesis.

For an additional exploratory diagnostic, assume each
$F_i(\pm i/2)=0$. Let $\psi=\Gamma'/\Gamma$,
let $\Lambda$ be the von Mangoldt function, and choose a finite
integer cutoff $N\geq2$. Define
\begin{align}\label{eq:finite-weil}
 K_{ij}
 &=\frac1\pi\int_0^\infty
   \operatorname{Re}\psi\!\left(\frac14+\frac{i\xi}{2}\right)
   F_i(\xi)F_j(\xi)\,d\xi-G_{ij}\log\pi,\\
 B^{(N)}_{ij}
 &=2\sum_{2\leq n\leq N}
   \frac{\Lambda(n)}{\sqrt n}\Phi_{ij}(\log n),
 \qquad W^{(N)}=K-B^{(N)}.\notag
\end{align}
The factor $2$ is the standard prime-side coefficient for
\eqref{eq:convention}; see \cite{conrey}. The software reports
numerical eigenvalues of $W^{(N)}$ as a \emph{finite-cutoff diagnostic}.
It supplies no bound on the omitted prime tail. Consequently,
neither the sign of $W^{(N)}$ nor that of a single numerical
quadratic-form value certifies the sign of the complete Weil form.

\section{Scope and reproducibility}

Equations \eqref{eq:inverse}, \eqref{eq:closed}, \eqref{eq:mass}, and
\eqref{eq:at-zero} are the formulas implemented by the accompanying
software. Numerical comparisons against independently integrated Fourier
inversion provide regression checks, not a substitute for the proofs above.
The polynomial formulas and $L^2$ Gram result above are algebraic and
unconditional. Finite Weil calculations use numerical quadrature and a
truncated prime-power sum, with no certified total error bound. The note
asserts no theorem conditional or unconditional on RH.

\begin{thebibliography}{9}
\bibitem{stein-shakarchi}
E. M. Stein and R. Shakarchi, \emph{Fourier Analysis: An Introduction},
Princeton University Press, 2003.
\bibitem{conrey}
J. B. Conrey, \emph{The Riemann Hypothesis},
Notices of the American Mathematical Society 50 (2003), 341--353.
\end{thebibliography}

\end{document}
